\section{Trabalhos Relacionados}

Este trabalho situa-se na interseção de duas literaturas: a de \textit{análise quantitativa do
mercado de transferências do futebol} e a de \textit{inferência causal com aprendizado de máquina}.
Esta seção revisa as frentes mais relevantes, destacando como cada uma fundamenta nossa metodologia
e em que pontos nossa abordagem se diferencia.

\subsection{Avaliação de valor de mercado de jogadores}

A estimativa do valor de um jogador é um problema central na economia do futebol. Uma linha
influente apoia-se na \textit{sabedoria das multidões}: as avaliações de mercado do Transfermarkt,
construídas a partir do julgamento coletivo de uma comunidade de usuários, mostram-se preditores
robustos de desfechos esportivos. Peeters \cite{peeters2018} demonstra que essas avaliações coletivas
predizem resultados de seleções nacionais com acurácia comparável ou superior à de \textit{rankings}
oficiais, evidenciando que o \textit{market value} carrega informação genuína sobre a qualidade do
atleta.

Nosso trabalho herda essa premissa, mas a aplica a uma tarefa distinta. Enquanto Peeters
\cite{peeters2018} valida o \textit{market value} contra resultados esportivos, nós o utilizamos como insumo de um
\textit{modelo hedônico} que estima o \textit{preço de transferência} (\textit{fee}) e, a partir do
resíduo desse modelo, isolamos o sobrepreço pago além do valor justo. Em vez de tomar o
\textit{market value} como verdade final, tratamo-lo como uma referência sobre a qual medimos
desvios conjunturais da negociação.

\subsection{O mercado de transferências como rede}

Uma segunda linha modela o mercado de transferências como uma \textit{rede dirigida e ponderada},
em que clubes são nós e transferências são arestas ponderadas pelo valor negociado. Esses estudos
caracterizam a estrutura do mercado — comunidades, centralidade e fluxos de capital — e relacionam a
posição estrutural de um clube ao seu desempenho.
% AMPLIAR SE HOUVER TEMPO: citar aqui \cite{wand2022,li2019,palazzo2023,dieles2023} (entradas
% comentadas em references.bib) ao descrever achados especificos da literatura de redes.

Diferentemente dessa abordagem, cujo foco está na \textit{estrutura} da rede, nosso trabalho
centra-se no \textit{preço} e no \textit{efeito causal} do histórico recente do comprador. Ainda
assim, a perspectiva de rede é incorporada de forma instrumental: métricas estruturais como
\textit{PageRank} e graus de entrada/saída entram em nosso modelo como \textit{confundidores}, para
neutralizar o fato de que clubes centrais na rede tendem a negociar valores maiores
independentemente de sua liquidez recente.

\subsection{Inferência causal com aprendizado de máquina}

Estimar um efeito causal a partir de dados observacionais exige separar a influência do tratamento
de confundidores que afetam tanto o tratamento quanto o desfecho. O arcabouço de \textit{Double/
Debiased Machine Learning} (DML) de Chernozhukov et al. \cite{chernozhukov2018} permite empregar modelos flexíveis
de aprendizado de máquina para estimar funções de confundimento de alta dimensão sem introduzir viés
de regularização no parâmetro causal, graças à ortogonalização de Neyman e ao \textit{cross-fitting}.
Adotamos esse arcabouço como núcleo da nossa estimação do efeito da liquidez sobre o sobrepreço.

Para investigar a \textit{heterogeneidade} do efeito — isto é, como ele varia entre clubes —
recorremos ao \textit{R-learner} de Nie e Wager \cite{nie2021}, que estima efeitos de tratamento
condicionais (CATE) por meio de uma formulação quasi-oráculo derivada da mesma lógica de
ortogonalização. Essa escolha nos permite construir um índice de vulnerabilidade por clube a partir
dos efeitos individuais estimados.

\subsection{Posicionamento}

Em síntese, combinamos as duas literaturas de uma forma ainda pouco explorada: aplicamos o
ferramental de inferência causal com aprendizado de máquina \cite{chernozhukov2018,nie2021} a uma
pergunta de economia do futebol tradicionalmente tratada de forma descritiva ou correlacional,
usando o \textit{market value} consagrado pela literatura \cite{peeters2018} como âncora para
definir o sobrepreço. O resultado é um desenho que distingue explicitamente o efeito causal da
liquidez do comprador do mero porte financeiro do clube — distinção que análises puramente
correlacionais ou estruturais não conseguem estabelecer.
